\cotefile{84}
\scanpage{82}


Donc la cat. des \unsure{-torseurs|torseurs} $2$-trivialisés, pour des Modules sur lesquels $2$ est inv., équivaut à celle des Modules sans plus.

On peut, sans \unsure{restr.} sur $2$ (\unsure{i.e.|inv.} sur $X$) ou sur $M$, regarder les \deleted{modules} torseurs $2$-trivialisés qui ont cette structure : munis de \inserted{scindages} $x_0\in\Gamma P_1$,
tels que $x_0=\overline{x_0}$, d'où $E\simeq L_0\oplus\underline{O}$, $T\simeq(0,2)$, comme dessus. On aura.

\[ b_{x_0,T}=2x_0-T=0,\]

La théorie est \unsure{tot.!|t.t.|tt. !} triviale \deleted{pour} $2$ inv. sur $X$ ! Le produit $*$ se réduit au produit $\otimes$ des \unsure{$L_i$|$L^{\mathrm i}$} sans plus, ``les $b$ suivent'' -- ils sont nuls !

Le cas \inserted{le plus \deleted{intéressant}} intéressant est celui où $2$ n'est \deleted{\unsure{ni nul,}} pas inv. sur $\mathcal{O}_X$\ldots

Formules

\[ \left\{\begin{aligned}\overline{x_1*x_2}&=\overline{x_1}*x_2=x_1*\overline{x_2}\\ \overline{x_1}*\overline{x_2}&=x_1*x_2\qquad\text{(résulte des précédentes)}\end{aligned}\right. \]

\deleted{i.e.} Prouvons $\overline{x_1}*x_2=\overline{x_1*x_2}$ il suffit pour $\unsure{x_2|x_1},x_2\in\Gamma P_{1}$

On a

\[ \begin{aligned}\overline{x_1}*x_2&=(T_1-x_1)*x_2=\underbrace{T_1}_{2x_1-b_1}*x_2-x_1*x_2=x_1*x_2-\\ &\hspace{18em}\deleted{\unsure{b}}b_1*x_2\\ \overline{x_1*x_2}&=T-x_1*x_2\deleted{\unsure{=}}=x_1*x_2-b_1\otimes b_2\end{aligned} \]

Or on a

\[ \begin{aligned}T&=2x_1*x_2-\underset{\substack{\Vert\\\deleted{2x_1-T_1}}}{b_1}\otimes\underset{\substack{\Vert\\\deleted{2x_2-\unsure{T_2|T_0}}}}{b_2}\qquad\text{d'où}\\ \overline{\unsure{a_1|x_1}*x_2}&=x_1*x_2-b_1\otimes b_2\qquad\text{OK.}\end{aligned} \]
% archive sheet number: 81
\scanpage{83}


\deleted{Supposons} Nous cherchons maintenant la condition pour que

a) $E_1\otimes E_2\longrightarrow E_1*E_1$ surjectif

b) Son noyau $N$ (qui contient les $\overline{\xi}-\xi$, $\xi\in E_1\otimes\overline{E_2}$) est engendré par \deleted{les} $\operatorname{Im}(\sigma_1\otimes\sigma_2-\mathrm{id})$, i.e. par les $\overline{x_1}\otimes\overline{x_2}-x_1\otimes x_2$ $(x_1\deleted{\unsure{,}}\in\Gamma E_1,\ x_2\in\Gamma E_2)\deleted{\unsure{x}}$, \deleted{\unsure{ou d'un|condition}} en d'autres termes on veut

\[ E_1\otimes E_2\mathbin{\unsure{-|\to}}E_1*E_2\quad\text{induit}\quad {E_1\otimes E_2}_{/\sigma_1\otimes\sigma_2}\simeq E_1*\unsure{I_2|E_2} \]

\deleted{\unsure{a}} Pour la condition a), elle équivaut à la surjectivité de

\[ L_1\otimes E_2\underset{L_1\otimes L_2}{\oplus}E_1\otimes L_2\longrightarrow L_1\otimes L_2,\]

or son conoyau est can. iso. à

\[ \operatorname{Coker}\pi_1\otimes\operatorname{Coker}\pi_2\qquad \pi_i:E_i\longrightarrow L_i \]

Si $L_1,L_2$ sont de présentation finie, alors la nullité (par Nakayama) se vérifie point par point, donc il est non seulement suffisant \inserted{(\unsure{tjs|après} OK)}, mais aussi néc., que \inserted{au voisinage de chaque point de $X$} $\pi_1$ \uuline{OU} $\pi_2$ soient épi. Si on a des scindages $x_1,x_2$,

d'où \quad $b_1\in\Gamma L_1,\ b_2\in\Gamma L_2$

alors

\[ \operatorname{Im}\deleted{\unsure{\pi_i:\,-}}(\pi_i:E_i\mathbin{\unsure{-|\to}}L_i)=\mathcal{O}_X\cdot b_i+2L_i. \]

donc la condition est la suivante :
% slanted

\margin{NB Cette \unsure{traduction|condition} N'utilise que la \unsure{condition|réduction} \unsure{sur|pour} $p=2$}

\textcircled{a$'$} \inserted{En tout} \deleted{Aux} point\deleted{s} de $X_0=V(\unsure{p|\gamma|2|\chi}\,.1)$, \deleted{\unsure{ou bien $b_1$ ou|ou $b_1$ ou $b_2$, ou}} \inserted{\unsure{Resp. ou|Resp. aux}} $T_{10}$, ou bien $T_{20}$, est un \deleted{\unsure{inv|base}} \inserted{gén.} de $L_{10}$ resp $L_{20}$.
% archive sheet number: 82
\scanpage{84}


Si on se donne (loc\textsuperscript{t.} ?) des bases \deleted{\unsure{de}} \inserted{$\{e_1\}$} de $L_1$, $\{e_2\}$ de $L_2$, donc $b_1,b_2$ s'interprèt\deleted{ent} comme des sections de $O$, alors la condition est

\[ (b_1,b_2,\underset{\substack{\Vert\\2\,\text{en}\\\text{l'occurrence}}}{\unsure{p|\gamma|\chi}})\cdot\underline{O}_X=\underline{O}_X \]

Cette condition est vérifiée si $p$ \inserted{(soit $2$ en l'occurrence)} inv. sur $X$. La condition b) aussi, comme on voit \deleted{en} la décomposition canonique

\[ E_i\simeq L_i\oplus O_i\qquad\text{avec }\sigma_i\mid O_i=+{\rm id}_{O_i} \]

d'où

\[ E_1\otimes E_2\simeq L_1\otimes L_2\oplus L_1\oplus L_2\oplus\underline O \]

et

\[ \sigma_1\otimes\sigma_2\quad\text{est}\quad\begin{cases}{\rm id}\text{ sur }L_1\otimes L_2\text{ et sur }O\\-{\rm id}\text{ sur }L_1\oplus L_2\end{cases} \]

et

\[ {\rm id}-\sigma_1\otimes\sigma_2\quad\text{est donc}\quad\begin{cases}0\text{ sur }L_1\otimes L_2\text{ et sur }O\\2\,{\rm id}\text{ sur }L_1\oplus L_2\end{cases} \]

donc

\diagram{\begin{tikzcd}[ampersand replacement=\&,column sep=small,row sep=large] (E_1\otimes E_2)_{/\sigma_1\otimes\sigma_2}\arrow[r,phantom,"\simeq"] \& (L_1\otimes L_2)\oplus O\arrow[r,"\sim"]\arrow[dr,"{x_1\otimes x_2\mapsto2x_1\otimes x_2}"',near end]\arrow[dr,bend left=18,"{\lambda\mapsto\lambda}"] \& \underbrace{E_1*E_2}_{\substack{\text{avec son}\\\text{scindage}\\\text{canonique}}}\\ \& \& L_1\otimes L_2\oplus O\end{tikzcd}}

Donc le seul pbl se pose en car. $2$, comme de juste. Je dis qu'\deleted{\unsure{alors}} alors (\unsure{Si on prend|Si on a} $L_1,L_2$ de prés. finie\unsure{\}|)}, p. ex. des Modules inversibles) la condition b) est \uline{conséquence} de a).
% slanted

\margin{\deleted{\unsure{On}} noter si on suppose que $L_1,L_2$ sont loc libres, alors dans $E=\unsure{\mathscr C_1|C_1}E_1\otimes E_2$ \unsure{$\mathscr C_i$} \unsure{$\operatorname{im}$} $(E_1\otimes E_2\longrightarrow E_1*E_2)$ la formation des $\unsure{\mathscr C_i}$ commute au changement de base $X_0\hookrightarrow X$.}

Il suffit le vérifier \deleted{\unsure{au ca}} si $X_{\!0}=\deleted{\unsure{X}}\unsure{X_0|X}$, de car. $2$ \textsuperscript{*} et il suffit de voir que si alors $T_1\in\Gamma L_1$
% archive sheet number: 83
\scanpage{85}


est inversible, alors le noyau de $*:E_1\otimes E_2\longrightarrow E_1*E_2$ est $\operatorname{Im}({\rm id}_{E_1\otimes E_2}-\sigma_1\otimes\sigma_2)$. [\uline{NB} Si $T_1$ et $T_2$ sont nuls, on a $\sigma_i={\rm id}$ $(i=1,2)$ donc $\mathrm{id}-\sigma_1\otimes\sigma_2=0$,
dont l'image est nulle, alors que le noyau de $E_1\otimes E_2\longrightarrow E_1*E_2$ \deleted{\unsure{contient}} \inserted{est ég. à} \inserted{celui de l'augmentation,} $E_1*_{0}E_2=L_1\otimes E_2+E_1\otimes L_2$ (qui est de rang $3$ si si $L_1,L_2$ inv.).]

En termes des scindages $x_1,x_2$ de $E_1,E_2$, et des \deleted{\unsure{la}} décompositions correspondantes

\[ \begin{aligned} E_1&\simeq L_1\oplus O\,x_1,\quad E_2\simeq L_2\oplus\underline O\,x_2\\ E_1\otimes E_2&\simeq(L_1\otimes L_2)\oplus\deleted{\unsure{\underline O}}\underset{\simeq L_1}{\underline{L_1\otimes x_2}}\oplus\underset{\simeq L_2}{\underline{x_1\otimes L_2}}+\underset{\simeq O}{\underline{\underline O\cdot x_1\otimes x_2}} \end{aligned} \]

et compte tenu que $\sigma_1\otimes\sigma_2\inserted{-{\rm id}}$ est \deleted{l'id.} \inserted{nul} sur le morceau $L_1\otimes L_2$, \deleted{il reste à voir} $\operatorname{Im}(\deleted{\unsure{{\rm id}-}}\sigma_1\otimes\sigma_2-{\rm id}_{E_1\otimes E_2})$ est engendré par $3$ sortes d'éléments, correspondants aux trois autres facteurs

\[ \begin{aligned} \overline{u_1\otimes x_2}-u_1\otimes x_2&=-u_1\otimes\underset{\substack{\Vert\\T_2-x_2\\\Vert\\x_2-b_2}}{\overline{x_2}}\unsure{-|=}u_1\otimes x_2=-\underline{2u_1}\otimes x_2+\underline{u_1\otimes b_2}\\ \overline{x_1\otimes u_2}-x_1\otimes u_2&=-2x_1\otimes\underline{u_2}+b_1\otimes u_2\\ \overline{x_1\otimes x_2}-x_1\otimes x_2&=\underset{(x_1-b_1)\otimes(x_2-b_2)}{\overline{x_1}\otimes\overline{x_2}}-x_1\otimes x_2=-x_1\otimes b_2-b_1\otimes x_2+b_1\otimes b_2 \end{aligned} \]

donc les éléments (avec les notations de la somme directe)

\[ \left\{\begin{aligned}&-2u_1+u_1\otimes b_2 &&u_1\in\unsure{T|\Gamma} L_1\\ &-2u_2+b_1\otimes u_2 &&u_2\in\Gamma L_2\\ &-b_1-b_2+b_1\otimes b_2\end{aligned}\right. \]

En car. $2$ cela se réduit à

\[ \left\{\begin{aligned}&u_1\otimes b_2,\ b_1\otimes u_2\in\Gamma L_1\otimes L_2\\ &-b_1-b_2-b_1\otimes b_2\quad\unsure{\notin}\end{aligned}\right. \]
% archive sheet number: 84
\scanpage{86}


Si $b_1$ \deleted{\unsure{en}} engendre $L_1$ (supposé inv.) alors les $b_1\otimes u_2$ engendrent $L_1\otimes L_2$, donc \deleted{\unsure{il}} $K=\operatorname{Im}(\sigma_1\otimes\sigma_2-{\rm id}_{E_1\otimes E_2})$ contient $L_1\otimes L_2$, plus précisément il est égal à

\[ L_1\otimes L_2+O\cdot\xi \]

où

\[ \xi=\overline{x_1\otimes x_2}-\unsure{x_2|x_1}\otimes x_2\equiv\deleted{b_2}b_1\oplus b_2\quad(\text{mod }L_1\otimes L_2) \]

\inserted{non nul mod. $L_1\otimes L_2$}

Si $L_1$ est de rang $1$, $L_2$ de rang $n$, $\operatorname{Im}(\sigma_1\otimes\sigma_2-{\rm id})$ est donc de rang \deleted{$1+n$,} égal à

\[ \operatorname{rang} L_1\otimes L_2+1=n+1 \]

tandis que

\[ \begin{aligned}\operatorname{rang}(\operatorname{Ker} *)&=\underset{\substack{\Vert\\\operatorname{rg} E_1\,\operatorname{rg} E_2\\\Vert\\2\cdot(n+1)}}{\operatorname{rang}(E_1\otimes E_2)}-\underset{\substack{\Vert\\1+\operatorname{rang}(L_1\otimes L_2)\\\Vert\\1+n}}{\operatorname{rang} E_1*E_2}\\ &=n+1\end{aligned} \]

d'où l'égalité.

On trouve donc

\uline{Proposition} : Supposons \deleted{\unsure{$E$}} $L_1,L_2$ de \inserted{loc. libres de} \deleted{type} fini.

Pour que $E_1\otimes E_2\mathbin{\unsure{-|\to}}E_1*E_2$ induise

\[ {E_1\otimes E_2}_{/\sigma_1\otimes\sigma_2}\xrightarrow{\sim}E_1*E_2, \]

il f. et s. que \deleted{\unsure{à son}} \deleted{tout point} \deleted{\unsure{s}} de $X_0=V(2.1_{\unsure{x|X}})$, $T_{10}$ \uline{ou} $T_{20}$ engendre $L_{10}$ resp. $L_{20}$. (Si $L_1,L_2$ sont tous les deux inversibles, et si $E_1,E_2$ munis de scindages, d'où $b_1\in\Gamma L_1,b_2\in\Gamma L_2$, cette condition équivaut à

\[ 2\,\mathcal O_X+b_1L_1^\vee+b_2 L_2^{\unsure{\vee|*}}=\underline O_X\quad\Bigr\} \]
% archive sheet number: 85
\scanpage{87}


\uline{Corollaire} C'est aussi la condition néc. et suff. pour que, pour tout Module $M$, on ait

\[ \underline{\operatorname{Hom}}_{\underline O}(E_1*E_2,M)\xrightarrow{\sim}\underline{\operatorname{Hom}}(E_1\otimes E_2,M)^{\sigma_1\otimes\sigma_2} \]

Posant $E_1^\vee=\unsure{\mathcal A_1|A_1}$, $E_2^\vee=\mathcal A_2$, $(E_1*E_2)^\vee\simeq\mathcal A$, cela implique que

\[ \mathcal A\longrightarrow\mathcal A_1\otimes\mathcal A_2 \]

induit un isomorphisme

\[ \mathcal A\xrightarrow{\sim}(\mathcal A_1\otimes\mathcal A_2)^{\sigma_1^\vee\otimes\sigma_2^\vee}, \]

et inversement si on veut que, ce \deleted{soit} reste vrai après tout changement de base.

\begin{center}$\backslash$\end{center}

Complément sur les $\unsure{\gamma|p|\chi}$ trivialisations, au dessus de $X_0=V(\gamma)$ i.e sur un schéma $X$ sur lequel $\gamma1_X=0$.

La donnée d'une $\gamma$-trivialisation pour un \inserted{$L$-}torseur $P$, comme $\gamma P=0.P\simeq L$, équivaut simplement à une section $T$ de $L$,

et $\pi:P\longrightarrow L$ est l'application constante $x\longmapsto\underset{\substack{\Vert\\0}}{\gamma x}-T=-T$.

\[ \pi(x)=-T\overset{\rm def}{=}b \]

\unsure{En tous cas|Un bon cas|Vu ! Dans le cas}

\[ E_1\otimes E_2\mathbin{\unsure{-|\to}}E_1*E_2 \]
% archive sheet number: 86
\scanpage{88}


est nul sur \deleted{\unsure{$E_2$|$E_1$}} $L_1\otimes L_2$ (car il y est égal à $\unsure{\gamma|p|\chi}\,{\rm id}_{L_1\otimes L_2}$), donc il se factorise par

\[ \underbrace{E_1\otimes E_2/L_1\otimes L_2}_{E_1\overset{\leftrightarrows}{\otimes}E_2}\longrightarrow E_1*E_2,\tag*{$(*)$} \]

or le premier mb s'interprète comme une extension

\diagram{\begin{tikzcd}[ampersand replacement=\&,column sep=small]0\arrow[r,no head] \& L_1\deleted{\otimes}\times L_2\arrow[r,no head] \& E_1\overset{\leftrightarrows}{\otimes}E_2\arrow[r,no head] \& \underline O\arrow[r] \&0\end{tikzcd}}

en fait, c'est l'extension qui correspond au torseur

\[ P_1\times P_2\qquad\text{sous }L_1\times L_2 \]

on peut aussi la décrire comme

\[ E_1\overset{\leftrightarrows}{\otimes}E_2\simeq E_1\underset{\underline O}{\times}E_2=\left\{\begin{aligned}&(x_1,x_2)\in E_1\times E_2\\ &\mid\psi_1(x_1)=\psi_2(x_2)\end{aligned}\right\} \]

On trouve alors un hom d'extensions

\diagram{\begin{tikzcd}[ampersand replacement=\&,column sep=small]0\arrow[r] \& L_1\times L_2\arrow[r]\arrow[d] \& \underset{=\operatorname{Env}(P_1\times P_2)}{E_1\overset{\leftrightarrows}{\otimes}E_2}\arrow[r]\arrow[d] \& \underline O\arrow[r]\arrow[d,equal] \&0\\0\arrow[r] \& L_1\otimes L_2\arrow[r] \& E_1*E_2\arrow[r,dash] \& \underline O\arrow[r,dash] \&0,\end{tikzcd}}

où

\[ L_1\unsure{*|\times}L_2\longrightarrow L_1\otimes L_2 \]

est donné par

\[ (u_1,u_2)\longrightarrow u_1\otimes b_2+b_1\otimes u_2 \]
% archive sheet number: 87
\scanpage{89}


On peut donc interpréter $E_1*E_2$ comme l'ext. de $\underline O$ par $L_1\otimes L_2$ déduite de l'ext. $E_1\overset{\leftrightarrows}{\otimes} E_2$ de $O$ par $L_1\times L_2$,
\unsure{ou ce qui est la même chose|cela est bien encore|via cet homom., et}, $P_1*P_2$ comme le torseur déduit du $(L_1\times L_2)$-torseur $P_1\times P_2$ par l'homomorphisme précédent.

Retour sur \unsure{les|le} \deleted{produit} rev. quadratique.\textsuperscript{*}
% slanted

\margin{* Dire d'abord \unsure{que, pour l'instant, sur une extension|ce que l'on entend par une extension} $0\longrightarrow L\longrightarrow E\longrightarrow\underline O\longrightarrow0$ $2$-trivialisée, avec $L$ inv., et une forme quadratique $Q$ ``admissible'' $\inserted{E\mathbin{-}L^{\otimes2}}$ (pour cette structure), -- i.e. a) $Q(u)=u^2$ si $u\in\Gamma L$ b) $T_Q=Q$, i.e. $Q$ de la forme $Q(x)=x^2-Tx\psi(x)+N\psi(x)^2$, où $N\in\operatorname{Sym}^2(E)$ convenable (\unsure{après} tel que $\forall x$, $Q(x)$ soit \unsure{d'év.} dans $L^{\otimes2}$). En termes d'un scindage de $E$, par \unsure{un|une} $x\in\Gamma P$, d'où $b=\pi_T(x)\in\Gamma L$, $Q$ doit être de la forme $Q(\underset{\substack{\rotatebox{-90}{$\in$}\\\Gamma O}}{\lambda}x+\underset{\substack{\rotatebox{-90}{$\in$}\\\Gamma L}}{u})=u^2+\lambda bu+\lambda^2c$ où $c\in\Gamma L^{\otimes2}$ -- le choix d'une $Q$ admissible revient donc à celui de $c\in\Gamma L^{\otimes2}$. \unsure{D'où} les $Q$ \unsure{admissibles} forment les sections d' un torseur sous $L^{\otimes2}$.}

Supposons \deleted{\unsure{à}} que $E_1,E_2$ soient associés à deux rev. quadratiques $X_1,X_2$ de $X$.

Cela signifie donc qu'on se donne sur $E_1,E_2$ formes quadratiques $Q_1,Q_2$

\[ \begin{aligned}Q_1:E_1&\mathbin{\unsure{-|\to}}L_1^{\otimes2}\\ Q_2:E_2&\mathbin{-}L_2^{\otimes2}\end{aligned} \]

dont les restrictions à $L_1,L_2$ sont $u\mapsto u^2$.

Ou encore que cela \unsure{r. à|R. à|$R_i$} des fonctions pol. de degré $\leq2$

\[ \begin{aligned}P_1&\longrightarrow L_1^{\otimes2},\\ P_2&\longrightarrow L_2^{\otimes2},\end{aligned} \]

dont les parties homogènes soient $u\longmapsto u^2$.

La donnée de $Q_1,Q_2$ définit des discriminants
% archive sheet number: 88
\scanpage{90}


\[ \delta_1\in\Gamma L_1^{\otimes\unsure{2|1}},\qquad\delta_2\in\Gamma L_{2}^{\otimes2} \]

donnés, en termes de scindages \inserted{(locaux ?)} $x_1,x_2$, par

\[ \delta_1=b_1^2-4c_1,\quad\delta_2=b_2^2-4c_2 \]

où

\[ b_1=\inserted{\pi_1(x_1)=}\underset{\substack{\rotatebox{-90}{$\in$}\\\Gamma L_1}}{2x_1-T_1},\qquad b_2=\pi_2(x_{\unsure{2|1}})=\underset{\substack{\rotatebox{-90}{$\in$}\\\Gamma L_2}}{2x_2-T_2} \]

et \quad $c_1=Q(x_1)\in\Gamma L_1^{\otimes2}$,\quad $c_2=\Gamma L_2^{\otimes2}$,

les formes quadratiques \inserted{sur $\deleted{\unsure{S}}E_1,E_2$} étant définies par

\[ \left\{\begin{aligned}Q_1(\lambda x_1+u_1)&=u_1^2\deleted{\unsure{-}}+\lambda b_1 u_1+\lambda^2 c_1\ ;\\ Q_2(\lambda x_2+u_2)&=u_2^2+\lambda b_2 u_2+\lambda^2c_2\end{aligned}\right. \]

donc sur $P_1,P_2$ par

\[ \begin{aligned}Q_1(x_1+u_1)&=u_1^2+b_1u_1+c_1\\Q_2(x_2+u_2)&=u_2^2+b_2u_2+c_2\end{aligned} \]

\uline{Th} Il existe une forme quadratique $Q$ ``admissible'' sur $E=E_1*E_2$ \inserted{à val. ds $L^{\otimes2}$ $(L=L_1\otimes L_2)$} (i.e. en telle que $Q\mid\deleted{\unsure{L_1}} L=L_1\otimes L_2$ \deleted{\unsure{à}} soit donnée par $u\longmapsto u^2$) et \inserted{\unsure{une}} seule, telle que l'on ait, pour

$x_1\in\Gamma \unsure{P_1|E_1}$, $x_2\in\Gamma \unsure{P_2|E_2}$ \deleted{(ou \unsure{plutôt} on suppose $x_1\in\Gamma P_1,x_2\in\Gamma P_2$)}

\[ \boxed{Q(x_1*x_2)=Q_1(x_1)\delta_2+Q_2(x_2)\delta_1+4Q_1(x_1)Q_2(x_2)} \]
% slanted

\margin{Si on suppose seulement $x_1\in\Gamma E_1$, $x_2\in\Gamma E_2$, on trouve $\boxed{\begin{aligned}\deleted{\unsure{Q_1(x_1)\delta_2}}Q(x_1*x_2)&=\psi_2(x_2)^2 Q_1(x_1)\delta_2+\psi_1(x_1)^2Q_2(x_2)\delta_1\\ &\quad+4Q_1(x_1)Q_2(x_2)\end{aligned}}$}

en termes des scindages définis par $x_1,x_2$, \inserted{$x_1*x_2$}, pour $x_1\in\Gamma P_1,x_2\in\Gamma P_2$

\[ \boxed{c=c_1\delta_2+c_2\delta_1+4c_1c_2} \]
% archive sheet number: 89
\scanpage{91}


i.e. 

\[ c=c_1(b_2^2-4c_2)+c_2(b_1^2-4c_1)+4c_1c_2\qquad\text{i.e.} \]

\[ \boxed{c=c_1b_2^2+c_2b_1^2-4c_1c_2}\ . \]

Comme on a de plus

\[ \boxed{b=b_1b_2}, \]

cela suffit à déterminer $Q$, comme

\[ Q(\underset{\substack{\rotatebox{-90}{$\in$}\\\Gamma O}}{\lambda}x_1*x_2\deleted{\unsure{)}}+\underset{\substack{\rotatebox{-90}{$\in$}\\\Gamma(L_1\otimes L_2)}}{u})=u^2+\lambda bu+\lambda^2c \]

\uline{Corollaire} Si $\delta$ est le discriminant de $Q$, on a

\[ \delta=\delta_1\delta_2\qquad\text{(en plus de }b=b_1b_2\text{)} \]

En effet

\[ \begin{aligned}\delta=b^2-4c&=b_1^2 b_2^2-4(c_1b_2^2+c_2b_1^2-4c_1c_2)\\ &=(b_1^2-4c_1)(b_2^2-4c_2)\end{aligned} \]

\uline{NB} Les formules simples à retenir, qui déterminent complètement $Q$ en termes de $Q_1$ et $Q_2$, quand on a des scindages, sont

\[ \boxed{\begin{aligned} b&=b_1b_2\\\delta&=\delta_1\delta_2\end{aligned}} \]

on en conclut pratiquement la formule pour $c$, puisque

\[ \delta=b^2-4c=\delta_1\delta_2\qquad\text{d'où, comme }b=b_1b_2 \]

\[ \begin{aligned}-4c=\delta_1\delta_2-b^2&=(b_1^2-4c_1)(b_2^2-4c_2)-b_1^2\cdot b_2^2\\ &=\unsure{4|4_1}(-b_1^2c_2\unsure{-|+}c_1b_2^2+4c_1c_2)\end{aligned} \]
% archive sheet number: 90
\scanpage{92}


i.e.

\[ 4c=4\left(b_1^2c_2+b_2^2c_1-4c_1c_2\right)=4(c_1\delta_2+c_2\delta_1+4c_1c_2) \]

\inserted{(\uline{\unsure{heuristiquement}})}

et la formule pertinente pour $c$ s'en déduit ``en divisant par 4''\ldots

Pour prouver le \unsure{th.|th}, il suffit de le faire localement, à cause de l'unicité, donc on peut choisir un scindage, et définir $Q$ \unsure{ad hoc|obtenue} par $b$ et $c$ donnés par les formules plus haut. Il faut prouver que $Q(y_1*y_2)$ satisfait la formule écrite -- c'est un petit calcul \unsure{délectable|détestable}, que j'ai pratiquement déjà fait plus haut.

\uline{Définition} de $X_1\unsure{\wedge|*|\star}X_2$ et de

\[ X_1\times X_2\longrightarrow X_1\wedge X_2 \]

Contraintes de comm, assoc, unité pour cette opération

$X_1\wedge X_2$ quand $X_1$ (ou $X_2$) est étale, comme déduit de $X_2$ en le ``tordant'' par le $(\unsure{\mathbb{Z}_{2\mathbb Z}|\mathbb Z/2\mathbb Z})_X$-torseur $X_1$

$X_1\wedge X_2$ est étale ssi $X_1$ et $X_2$ le sont

$X_1\times X_2\unsure{-|\to}X_1\wedge X_2$ est fid. plat, sauf
% archive sheet number: 91
\scanpage{93}


exactement en les points de $X$ \deleted{$\unsure{x\in V(2\cdot1_X)|x\in V(2\,1_X)}$} \unsure{(lire|lire} \inserted{ds $X$}\unsure{)|\,} \deleted{de car. 2} où $X_1$ \uline{et} $X_2$ sont non étales.

\deleted{Calculs en termes des scindages, avec $x_1\in\unsure{\Gamma P_1|\Gamma xP_1},\ x_2\in\Gamma P_2$, d'où}

\[ \deleted{E_1\simeq L_1\oplus\underset{\substack{\rotatebox{90}{$=$}\\O.x_1}}{O},\qquad E_2\simeq L_2\oplus\underset{\substack{\rotatebox{90}{$=$}\\O.x_2}}{O}} \]

\deleted{et}

\[ \deleted{\mathcal{A}_1=O\unsure{\oplus|+}L_1^\vee,\qquad\mathcal{A}_2\simeq\underline O+L_2^\vee} \]

\deleted{\deleted{\unsure{Pour $d_1$, la forme $x$|Donc $d_1$, la forme $x$|Donc $L_1$, les formes $x$}}}

\deleted{On définit $L_1,L_2$ commes sous-mod. de $E_1,E_2$ \unsure{par|Posons}}

\[ \deleted{L_1=\operatorname{Ker}(x_1:\mathcal{A}_1\unsure{-|\to}O),\quad L_2=\operatorname{Ker}(x_2:\mathcal{A}_2-O)} \]

\deleted{On a}

\[ \deleted{b_1=2x_1-T_1\qquad\text{i.e.}\quad \deleted{b_1(u'_1)}\inserted{b_1.u'_1}=\deleted{\unsure{+|-}}-\operatorname{Tr}(u'_1)} \]

\deleted{et on a}

\[ \deleted{N(u'_1)={u'_1}^2.c_1} \]

\deleted{donc pour $u'_1\in L_1^\vee$, l'équation universelle \inserted{dans $\mathcal{A}_1$,} est}

\[ \deleted{{u'_1}^2+b_1 u'_1+c_1=0} \]

\deleted{On a}

\[ \deleted{N(\overbrace{u'_1+\underset{\substack{\rotatebox{90}{$=$}\\\deleted{\unsure{u_1}}\\x_1(\xi_1)}}{\lambda}e_1}^{\xi_1})=\deleted{\unsure{{u'}^2.1+\lambda u'+\lambda|{u'}^2.1+\lambda^2+\lambda}}\inserted{\underset{={u'}^2.c_1}{N(u')}+\underset{\substack{\rotatebox{90}{$=$}\\\lambda u'.b_1}}{\lambda\operatorname{Tr}(u')}+\underset{\substack{\rotatebox{90}{$=$}\\x_1^2(\xi_1)}}{\lambda^2}}} \]

\deleted{donc}

\[ \deleted{\deleted{\unsure{N-x_1|N=x_1}}=\deleted{\unsure{N}}{u'}^2.c_1+x_1(\xi_1} \]
% archive sheet number: 92
\scanpage{94}


Retour sur \deleted{\unsure{leurs|leur|une}} \unsure{rev.|rev} quadratique, $X'/X$ défini par $\mathcal{A}$ alg. quadratique sur $\overline O$, ou par $(E,\psi:E\unsure{-|\to}O,T,Q)$. Choisissons un scindage $x_0\in\Gamma P_{1}=\psi^{-1}(1)$, alors

\[ E\simeq L\oplus\underset{\substack{\rotatebox{90}{$\simeq$}\\O}}{O.x_0},\qquad \mathcal{A}\simeq O\oplus L^\vee \]

où $L^\vee$ est défini comme sous-module de $\mathcal A$ \deleted{\unsure{par}} comme $\operatorname{Ker}(x_0:\mathcal{A}\to\underline O)$.

Considérons

\[ \boxed{\underset{\substack{\rotatebox{-90}{$\in$}\\\Gamma L}}{b}=2x_0-T} \]

donc pour $u\in L^\vee$, on a

\[ \langle u,b\rangle\overset{\mathrm{def}}{=}\boxed{u.b=-\operatorname{Tr}(u)} \]

On aura

\[ N=x_0^2\deleted{\unsure{+}}-bx_0+c\in\operatorname{Sym}^2(E),\quad\inserted{\text{où }}c\in\Gamma L^{\otimes2} \]

donc pour $u\in L^\vee$, on aura

\[ N(u)=u^2.c\qquad(\overset{\mathrm{def}}{=}\langle u^{\otimes2},c\rangle) \]

donc l'équation minimale de $u$ \deleted{dans} $\mathcal{A}$ est

\[ u^2+\underbrace{\langle u,b\rangle}_{\substack{\rotatebox{90}{$=$}\\ b(u)\\\deleted{\unsure{u}}}}\,\unsure{u|\,}+\underbrace{\langle u^{\otimes2},c\rangle}_{\substack{\rotatebox{90}{$=$}\\c(u)}}=0 \]

\deleted{donc on aura}

On a, pour $u\in\Gamma L^\vee$, $\lambda\in\Gamma O$

\[ \begin{aligned}N(\underbrace{u+\lambda}_{\xi})&=N(u)+\lambda\operatorname{Tr}(u)+\lambda^2\qquad\text{où }\lambda=x_0(\xi)\\ &=x_0(\xi)^2\deleted{\unsure{+}}-x_0(\xi)\underset{\substack{\rotatebox{90}{$=$}\\b(\xi)}}{b(u)}+\underset{\substack{\rotatebox{90}{$=$}\\c(\xi)}}{c(u)}\end{aligned} \]

donc on aura bien

\[ N=x_0^2-x_0 b+c \]

D'autre part, si $x=x_0+v\in\Gamma P$ \quad $(v\in\Gamma L)$ \deleted{\unsure{on aura}} $x$ sera donc la \deleted{\unsure{fonction|application}} \inserted{$\underline O$-linéaire $\mathcal A\to\underline O$} \inserted{$\xi:\unsure{\xi|\,}=u+\lambda\longmapsto\lambda+\langle u,v\rangle$} qui sur $1$ est \unsure{$=1$|$^21$|$1$}, et sur $L$ \inserted{$u\longmapsto\langle u,v\rangle$} est donné par
% archive sheet number: 93
\scanpage{95}


et on vérifie tout de suite que c'est un hom d'alg. ssi on a

\[ \underbrace{v^2+bv+c}_{\substack{\rotatebox{-90}{$\in$}\\\Gamma L^{\otimes2}}}=0\qquad\text{i.e.}\quad Q(v)=0 \]

On va préciser la relation entre les formes $N$ sur $\mathcal A$, $Q$ sur $E=\mathcal A^\vee$, \deleted{par la forme} \inserted{en utilisant} l'iso (valable pour \unsure{tt|tout} $E$ loc libre de rang 2)

\[ E\xrightarrow{\sim}\underline{\operatorname{Hom}}(E,\underbrace{\det E}_{\substack{\rotatebox{90}{$=$}\mathrm{def}\\\Delta}})\simeq\underset{\substack{\rotatebox{90}{$=$}\mathrm{def}\\\mathcal A}}{E^\vee}\otimes\Delta=\mathcal A\otimes\Delta \]

obtenu en associant à $x\in E$ la forme

\[ f_x:E\unsure{-|\to}\det E=\mathop{\wedge}\nolimits^2 E\qquad y\longrightarrow x\wedge y \]

Si $E$ est muni d'une structure d'extension

\diagram{\begin{tikzcd}[ampersand replacement=\&]0\arrow[r,no head]\&L\arrow[r]\&E\arrow[r,"\psi",no head]\&O\arrow[r]\&0,\end{tikzcd}}

alors

\[ \underset{\substack{\rotatebox{90}{$=$}\\\Delta}}{\det E}\xleftarrow[\sim]{\alpha}L\qquad\text{(avec conv. de signe à préciser)} \]

\[ \alpha(u)=x_0\wedge\unsure{u|v},\quad\text{où }x_0\in\Gamma P_1 \]

\margin{\uline{NB} le $2^{\mathrm e}$ mb ne dépend pas du choix de $x$}

donc

\[ E\xrightarrow{\sim}\underset{\substack{\rotatebox{90}{$=$}\\E^\vee}}{\mathcal A}\otimes L \]

À tout $x\in E$ on associe la forme

\[ g_x:E\longrightarrow L \]

\unsure{donné|donnée} par la condition

\[ g_x(y)=\alpha^{-1}(x\wedge y)\ \deleted{\unsure{x}} \qquad\text{i.e. par la condition} \]

\[ x_0\wedge g_x(y)\deleted{\unsure{x_0|u}}=x\wedge y \]

\deleted{Si $T\in E$ tel que $\psi(T)=2$,} je dis qu'on aura \deleted{pour \unsure{prendre|prouver}}

\[ g_x(y)=\psi(x)y-\psi(y)x \]
% archive sheet number: 94
\scanpage{96}


Pour le vérifier, on note que, pour $x,y$ variables, les deux mbs de la formule sont des formes bil. alternées, et $\{(x,y)\unsure{-|\mapsto}g_x(y)=\alpha^{-1}(x\wedge y)\}$ forme une base du \unsure{Module|module} \unsure{(d'inv.)|(inv.)} des telles formes alternées. Donc on doit avoir

\[ \psi(x)y-\psi(y)x=\unsure{\lambda'|\lambda}\,\alpha^{-1}(x\wedge y)\qquad\text{i.e.} \]

\deleted{avec}

\[ \alpha(\psi(x)y-\psi(y)x)=\lambda\,x\wedge y, \]

\deleted{\unsure{pour}} \unsure{où|\,} $\lambda\in\Gamma O$ est une section \unsure{indépend.|indépend.te} de $x,y$. On veut prouver $\lambda=1$, il suffit donc de le vérifier loc, pour des $x,y$ tels que $x\wedge y$ soit une base de $\det E$. Soit $u$ base de $L$, $x_0$ tel que $\unsure{\varphi|\psi|\Psi}(x_0)=1$, on aura (en prenant $x=u$, $y=x_0$)

\[ \alpha(-\unsure{u|x})=\lambda\underbrace{u\wedge x_0}_{-\alpha(u)} \]

\[ \text{i.e.}\qquad u=\lambda u \]

donc \deleted{(manque de pot !)} $\lambda=1$.

Notons que

\diagram{\begin{tikzcd}[ampersand replacement=\&]0\arrow[r,no head]\&\unsure{O|\mathcal A}\arrow[r,no head]\&\mathcal A\arrow[r,"\varphi",no head]\&L^\vee\arrow[r,no head]\&0\end{tikzcd}}

(transposée de $0-L-E-\underline O-0$)

donne par tensorisation

\diagram{\begin{tikzcd}[ampersand replacement=\&]0\arrow[r,no head]\&L\arrow[r,no head]\&\mathcal A\otimes L\arrow[r,"\psi",no head]\&\underline O\arrow[r,no head]\&0\end{tikzcd}}

et je dis que $g:E-\mathcal A\otimes L$ est un hom d'extensions
% archive sheet number: 95
